% =====================================================================
% kth-style.tex - KTH-flavored typographical conventions and a title-page
% macro mimicking the KTH EECS master thesis layout. This is NOT the
% official KTH template - it is a lightweight stand-in that follows the
% same visual conventions (Latin Modern, A4, centered title page with
% school + program + TRITA fields).
% =====================================================================

% ---- Headers / footers ----------------------------------------------
\usepackage{fancyhdr}
\pagestyle{fancy}
\fancyhf{}
\renewcommand{\chaptermark}[1]{\markboth{\chaptername\ \thechapter.\ #1}{}}
\fancyhead[L]{\nouppercase{\leftmark}}
\fancyhead[R]{\thepage}
\renewcommand{\headrulewidth}{0.4pt}

\fancypagestyle{plain}{%
  \fancyhf{}%
  \fancyfoot[C]{\thepage}%
  \renewcommand{\headrulewidth}{0pt}%
}

% ---- Chapter heading style (compact, KTH-ish) -----------------------
\usepackage{titlesec}
\titleformat{\chapter}[display]
  {\normalfont\sffamily\Large\bfseries}
  {\MakeUppercase{\chaptertitlename}\ \thechapter}
  {0.5em}
  {\Huge}
\titlespacing*{\chapter}{0pt}{-20pt}{30pt}

% ---- Thesis metadata (override in main.tex if desired) --------------
\newcommand{\thesistitle}{Joint Radar Emitter and Mode Clustering Using a Transformer Encoder Architecture}
\newcommand{\thesissubtitle}{}
\newcommand{\thesisauthor}{Markus Swegmark}
\newcommand{\thesisyear}{\the\year}
\newcommand{\thesismonth}{\ifcase\month\or January\or February\or March\or April\or May\or June\or July\or August\or September\or October\or November\or December\fi}

\newcommand{\thesisschool}{School of Electrical Engineering and Computer Science}
\newcommand{\thesisuniversity}{KTH Royal Institute of Technology}
\newcommand{\thesisprogram}{Master's Programme in Machine Learning}
\newcommand{\thesissupervisor}{TBD}
\newcommand{\thesisexaminer}{TBD}
\newcommand{\thesishost}{TBD}
\newcommand{\thesistrita}{TRITA-EECS-EX-XXXX:XXX}

\hypersetup{
  pdftitle={\thesistitle},
  pdfauthor={\thesisauthor},
  pdfsubject={Master's Thesis, KTH EECS},
  pdfkeywords={transformer, clustering, radar, deep learning, ELINT}
}

% ---- Title-page macro ----------------------------------------------
% Renders a centered KTH-style title page. Logos can be dropped into
% figures/kth-logo.pdf later; we wrap the include in \IfFileExists so
% the build does not fail when the asset is missing.
\newcommand{\maketitlepageKTH}{%
  \begin{titlepage}
    % Suppress hyperref's "page.N" destination so the title page does
    % not collide with later page-1's (abstract, first chapter).
    \hypersetup{pageanchor=false}%
    \centering
    \IfFileExists{figures/kth-logo.pdf}{%
      \includegraphics[width=0.25\textwidth]{figures/kth-logo.pdf}\par
    }{%
      {\sffamily\bfseries\Large KTH}\par
    }
    \vspace{1.5cm}
    {\sffamily\large \MakeUppercase{Degree Project in Computer Science and Engineering}\par}
    {\sffamily\large \MakeUppercase{Second Cycle, 30 Credits}\par}
    \vspace{0.5cm}
    {\sffamily\small Stockholm, Sweden \thesisyear\par}

    \vfill
    {\Huge\bfseries\thesistitle\par}
    \ifx\thesissubtitle\empty\else
      \vspace{0.4cm}{\Large\itshape\thesissubtitle\par}
    \fi
    \vspace{1.5cm}
    {\Large\thesisauthor\par}

    \vfill
    \begin{tabular}{rl}
      Supervisor: & \thesissupervisor \\
      Examiner:   & \thesisexaminer   \\
      Host:       & \thesishost       \\
    \end{tabular}\par

    \vspace{1cm}
    {\sffamily\small \thesisuniversity\par}
    {\sffamily\small \thesisschool\par}
    {\sffamily\small \thesisprogram\par}
    \vspace{0.5cm}
    {\sffamily\footnotesize \thesistrita\par}
  \end{titlepage}
  \hypersetup{pageanchor=true}%
}
